\section{Proposed System: A Self-Evolving Omni Speech Agent}
\label{sec:system}

We now instantiate the theory of \Cref{sec:osa} as a concrete, buildable agent. The design is governed by a single discipline: every component that \emph{evolves} at inference time does so only through an acceptance gate driven by a reward signal and operating inside a $\beta$-KL trust region, never through a gradient on the base model. This is what distinguishes our agent from the prevailing self-evolving paradigm, whose append-only memories and self-generated skill libraries are documented to drift toward context collapse and error propagation \citep{ace2025,expfollowing2025,remembering2026}. We design \emph{against} that failure mode by realizing the optimization geometry of $\qs$ (the Gibbs optimum of \Cref{sec:osa}) as an external, mutable state whose update operator is \emph{engineered to be contractive}. We stress at the outset that contractivity here is a design target, not a theorem: as \Cref{sec:convergence} makes explicit, finite-time monotone-improvement over the full agent state is the open problem this work frames, and the design principles below are our attempt to approach it---not a certificate that they attain it.

\subsection{Architecture}
\label{subsec:arch}

The agent couples two \emph{frozen} omni models with one external evolving state. A frozen \emph{vector-omni} encoder \citep{omniembed2025,omniembedaudio2026} serves as the memory and skill index: it maps raw spoken turns into the embedding action space $\cZA$ of \Cref{sec:osa}, where retrieval is a tilting of the base index distribution $\qo$. A frozen \emph{generative omni} policy \citep{qwen3omni} serves as the response model, whose decoding occupies the generative action space $\cZB$. Neither model is fine-tuned; both are queried only. All adaptation lives in an external state $S$---a structured memory store and a skill library---which is mutated under a mixture of \emph{verifiable} speech rewards and \emph{surrogate} self-consistency signals (detailed below). Formally, the agent's inference-time degrees of freedom $\cZ = \cZA \times \cZB$ are conditioned on $S$, and learning is the trajectory $S_0 \to S_1 \to \cdots$ of \emph{state} updates rather than weight updates. Because $S$ is the only mutable object, the convergence question reduces to the contractivity of the state-update operator---which we engineer directly, while acknowledging (\Cref{sec:convergence}) that the resulting trajectory is at best asymptotically controllable and is not yet certified to converge in finite time.

\paragraph{Which rewards are verifiable.} We hold to the paper's definition: a reward is \emph{verifiable} only if it is computable from a ground-truth annotation without a learned scorer. Under this definition the agent's genuinely verifiable rewards are: \textbf{WER/CER} for ASR (against reference transcripts), \textbf{exact-match} for intent/slot tasks (against gold labels), and \textbf{speaker-ID accuracy} \emph{only on corpora that carry ground-truth speaker labels}. Paralinguistic agreement signals (speaker/emotion consistency against stored estimates) do \emph{not} meet this bar and are treated separately as surrogates (\Cref{subsec:memory}). The architecture foregrounds the verifiable channel: wherever a ground-truth annotation exists, the acceptance gate uses it; surrogates are used only to fill gaps and are explicitly flagged as Goodhart-prone.

\paragraph{Why two frozen models.} Separating the index ($\cZA$) from the policy ($\cZB$) lets each operator be tilted by an independent reward channel: memory writes are gated by retrieval-utility signals, while response decoding is gated by task rewards. This factorization mirrors the Operator-A / Operator-B decomposition of \Cref{sec:osa} and lets us bound the two mutation rates separately (\Cref{subsec:control}).

\subsection{Memory}
\label{subsec:memory}

\paragraph{Three tiers over a bi-temporal graph.} Memory is a structured store with three tiers: \emph{episodic} (individual spoken turns), \emph{semantic} (per-speaker persona and stable facts), and \emph{procedural} (pointers into the skill library of \Cref{subsec:skills}). The tiers sit on a bi-temporal knowledge graph that records both event time and ingestion time \citep{arigraph2024,zep2025,amem2025}, so that retrieval can respect causality and the curator can reason about staleness. Episodic nodes consolidate into semantic nodes off the critical path.

\paragraph{Index on audio, not transcripts.} A central design choice is that the retrieval key is the frozen \emph{audio} embedding from the vector-omni encoder, not an ASR transcript \citep{wavrag2025,voxrag2025}. Transcribing before indexing discards exactly the paralinguistic structure (prosody, speaker identity, affect) that the agent must remember, and it injects ASR error into the key. Indexing on $\cZA$ keeps the key in the acoustic space.

\paragraph{Decoupled retrieval keys: do not overload the content embedding.} A caveat is essential here. Our own probing results (see the feasibility tables) indicate that the frozen vector-omni \emph{content} embedding is at or near chance on speaker identity and only weakly informative for emotion---it is optimized for lexical/semantic content, not paralinguistics. We therefore do \emph{not} route paralinguistic retrieval through the content embedding. Instead the composite key is built from \emph{dedicated} representations: a speaker subkey from frozen ECAPA-TDNN / x-vector embeddings \citep{ecapatdnn2020,xvector2021}, an emotion subkey from a dedicated SER embedding, and a content subkey from the omni content embedding. Paralinguistic queries retrieve on the speaker and SER subspaces; content queries retrieve on the content subspace. The content embedding is reserved for what it is good at. This decoupling is a design recommendation conditioned on the probing evidence; the exact gating weights are an empirical question we leave to evaluation.

\paragraph{Composite keys and their reward status.} Each memory entry carries a composite key
\[
k = (\text{speaker-ID},\ \text{emotion/SER},\ \text{turn-index}),
\]
extracted by frozen diarization, an x-vector / ECAPA-TDNN speaker embedding, and an SER head \citep{xvector2021,ecapatdnn2020}. These fields are \emph{dual-use} as retrieval structure and as a feedback channel---but we are careful about what kind of feedback. When the agent later answers a speaker- or emotion-conditioned query, agreement between a freshly predicted attribute and the stored key field is \textbf{not} a verifiable reward in the sense of this paper: the stored key is itself a frozen-model estimate (diarization / x-vector / ECAPA / SER output), so ``agreement'' measures \emph{self-consistency} of two model passes, not consistency with ground truth. We therefore reclassify key-agreement as a \emph{surrogate self-consistency signal}. It is useful---it is cheap, label-free, and correlates with stability---but it is Goodhart-prone: a frozen extractor's systematic error becomes a fixed point the agent can satisfy without being correct, so optimizing it can reinforce the extractor's bias rather than the truth. We use it only as an auxiliary signal, down-weighted relative to verifiable rewards, and never as the sole admission criterion for a write. Where corpus speaker labels exist, we substitute true speaker-ID accuracy for the surrogate. This is a deliberate departure from approaches that close the loop with a \emph{trained} memory manager \citep{memoryr1_2025}: we obtain a weak, honestly-labeled self-consistency signal from the keys, not a verifiable reward, and we lean on the genuinely verifiable channels (WER, exact-match, labeled speaker-ID) for the load-bearing credit assignment.

\paragraph{Retrieve--rerank.} Retrieval is two-stage. A fast first stage does late-interaction matching in $\cZA$ \citep{colbert2020}; a precise second stage reranks the shortlist by reward-weighted importance, weighting each candidate by its historical contribution to verifiable reward (with surrogate signals entering only as a discounted secondary term) \citep{genagents2023}. The reranker is the discrete analogue of the importance weight $\exp(\Rwd/\beta)$ that defines $\qs$ in \Cref{sec:osa}.

\paragraph{Non-destructive curation with provenance.} Writes are mediated by a curator that emits one of $\{\textsc{Add}, \textsc{Update}, \textsc{Delete}, \textsc{Noop}\}$ \citep{mem0_2025}; \textsc{Noop} is first-class and dominant, which is what keeps the store from bloating. Salience decays over time on an Ebbinghaus schedule \citep{memorybank2023}, so unreinforced episodic traces fade unless consolidated. Every node carries provenance and a trust score; low-trust or anomalous writes are quarantined to resist memory poisoning \citep{agentpoison2024}, the dominant attack surface for evolving stores.

\paragraph{Gradient-free credit assignment.} The curator does \emph{not} learn its policy by gradient. Given a candidate set of memory-write operations, it admits the write only if a best-of-$N$ comparison shows a positive reward delta against the no-write baseline---computed on the verifiable channel wherever ground truth is available, and otherwise on the discounted surrogate. This replaces the trained manager of \citet{memoryr1_2025} with a training-free selection rule whose acceptance weight is the tilting weight of \Cref{sec:osa}, inheriting the best-of-$N$ KL bounds of \citet{beirami2024bon,verdun2025softbon}---noting, as \Cref{sec:convergence} records, that those bounds are finite-$N$ guarantees for a \emph{fixed} operator and transfer only conditionally once the operator (the store) is itself mutated.

\paragraph{$\beta$-KL trust region on mutation rate.} The first instantiation of the control law bounds the \emph{rate} of memory mutation: the curator may not move the store's distribution over keys by more than a $\beta$-controlled KL budget per episode,
\[
\KL{q_{S_{t+1}}}{q_{S_t}} \le \epsilon_{\mathrm{mem}}(\beta).
\]
This trust region is the mechanism by which we \emph{design against} the context collapse documented for append-only agents \citep{ace2025,expfollowing2025}: a bounded-KL update operator is intended to be contractive in the sense \Cref{sec:osa} requires, so that the state trajectory is steered away from the divergent feedback loop rather than left to drift. We state this as a design principle, not a theorem: bounding per-step KL is necessary for stability but we do not claim it is sufficient for finite-time convergence, which remains open (\Cref{subsec:gap}). Consolidation (episodic$\to$semantic compaction, decay sweeps, re-embedding) runs asynchronously off the critical path, in the spirit of sleep-time and asynchronous-retrieval memory services \citep{lettasleeptime,moshirag2026}, so the trust region never blocks live inference.

\subsection{Skills}
\label{subsec:skills}

\paragraph{The speech-skill object.} A skill is a structured \textsc{Skill.md}-style object $(C, \pi, T, R)$---a usage context $C$, a callable procedure $\pi$, a set of test cases $T$, and a verifiable reward $R$ \citep{sokskills2026,anthropicskills}. The procedure $\pi$ is drawn from a small set of speech-native bodies: generative error correction for ASR rare-word repair \citep{gerrarewords2025}, voice tool-use orchestration \citep{aura2025}, and paralinguistic-preserving speech-to-speech rewriting \citep{paras2s2025}. Each body is a callable over the frozen policy, never a fine-tune of it.

\paragraph{Verification gate.} A candidate skill is admitted to the library only if it produces a \emph{positive verifiable-reward delta} relative to the no-skill baseline on its test set $T$; a resident skill whose running delta turns negative is deprecated \citep{skillsbench2026}. Here we require the delta on a genuinely verifiable channel (WER, exact-match, or labeled speaker-ID), precisely because the skill gate is load-bearing and must not be satisfiable by a surrogate. This operationalizes the curation-quality result of the skills literature: curated, verified skills dominate self-generated ones, and the gate is what enforces curation. At selection time the agent retrieves skills body-first---matching the situation to skill \emph{bodies} rather than to their natural-language descriptions \citep{skillrouter2026}---which avoids the description-mismatch failure of name-based routing.

\paragraph{Gradient-free improvement under a Pareto trust region.} Skills improve without touching base weights through reflective, evolutionary search: GEPA-style Pareto editing of skill prompts and procedures \citep{gepa2025} and AFlow-style workflow search over skill compositions \citep{aflow2024}. The governing \emph{design principle} is the second instantiation of our control law: an edited skill is accepted only if it is \emph{Pareto non-regressive}---it must not lose verifiable reward on any task in its competency set. We intend Pareto non-regression to act like a $\beta$-KL trust region on the skill library's behavior, keeping the post-edit policy within a bounded-KL neighborhood of the pre-edit policy on each verified task, so as to \emph{discourage} the reflection-plateau and skill-thrash pathologies \citep{reflexion2023,lats2023}. We do \emph{not} claim this guarantees monotone non-degradation: non-regression is checked only on the finite, observed competency set $T$, so it bounds measured regression on $T$ but cannot certify global non-degradation, and the convergence of the resulting edit sequence is exactly the open problem of \Cref{subsec:gap}. The construction is engineered to be contractive; certifying that it is remains future work.

\paragraph{Bloat control and composition.} Library growth is contained by polymorphic abstraction---collapsing near-duplicate skills into parameterized families \citep{polyskill2025}---and by CRAFT-style validate-and-deduplicate at admission \citep{craft2023}. Genuinely new competence accretes by AWM-style snowballing, where verified skills compose into higher-order workflows that are themselves gated \citep{awm2024}. Thus the library grows in capability while its effective description length stays bounded, the skill-side analogue of the memory decay schedule.

\subsection{The Unified Control Law}
\label{subsec:control}

The two subsystems are governed by one law, stated once:
\begin{quote}
\emph{Admit an inference-time state mutation iff it yields a positive reward delta against the no-mutation baseline---measured on a verifiable channel wherever ground truth exists, and otherwise on a down-weighted surrogate (the acceptance gate)---and only within a $\beta$-KL trust region that bounds how far the induced action distribution may move (the stability constraint).}
\end{quote}
The acceptance gate is the discrete realization of the reward tilting that defines the Gibbs optimum $\qs \propto \qo\exp(\Rwd/\beta)$ of the objective $\Fobj{q} = \Exp_q[\Rwd] - \beta\,\KL{q}{\qo}$; at the output level it inherits the $O(1/N)$ tilting rate and the KL--reward trade-off of \Cref{sec:convergence}. The trust region is the regularizer $\beta\,\KL{\cdot}{\qo}$ that keeps the search distribution $\qz$ from collapsing. The law has exactly two instantiations: the memory-mutation-rate bound $\epsilon_{\mathrm{mem}}(\beta)$ of \Cref{subsec:memory}, and the skill-library Pareto non-regression bound of \Cref{subsec:skills}. Both are bounded-KL, and both are \emph{engineered} to yield contractive state-update operators. We are deliberate about the strength of this claim: the finite-$N$ convergence results of \Cref{sec:convergence} hold for the \emph{fixed}-operator (output-level) case, whereas here the operator---the external state $S$---is itself part of the action, which is exactly the regime where only asymptotic, trust-region-conditioned guarantees are currently available and finite-time monotone convergence is unproven (\Cref{subsec:gap}). We therefore present the control law as the architecture's \emph{convergence design principle}: each admitted step is constructed to strictly improve a (preferably verifiable) reward while staying inside a $\beta$-controlled neighborhood, with no base-model gradient anywhere in the loop. Whether the resulting trajectory $S_0 \to S_1 \to \cdots$ provably converges---rather than merely being steered toward convergence---is the central open question this proposal sets out to answer, not a result we assert here.